\section{Minimum Number of Platforms Required for a Railway}
\label{sec:minimum-platforms}

\problemstatement{We are given the arrival times $\textit{arr}[0..n-1]$ and
departure times $\textit{dep}[0..n-1]$ of $n$ trains scheduled at a single
railway station, where train $i$ occupies a platform from
$\textit{arr}[i]$ until $\textit{dep}[i]$ (inclusive; if one train departs
at the exact same time another arrives, they are considered to overlap and
require separate platforms). We must find the \textbf{minimum number of
platforms} the station needs so that no train is ever forced to wait for a
free platform.}

\begin{patternbox}
\emph{``Minimum number of resources (platforms/rooms/servers) needed so that
no request ever waits, given many overlapping time intervals''} is the
classic \emph{maximum overlap / interval point sweep} pattern -- it is
closely related to, but distinct from, the activity-selection pattern in
Section~\ref{sec:n-meetings}. There, we had \emph{one} room and asked how
many of the requests we could serve; here, every request \emph{must} be
served, and we ask how many rooms that requires. The keyword shift from
``maximum number you can fit in one resource'' to ``minimum resources needed
to fit everyone'' is the signal to switch from finish-time-only greedy
scheduling to a \emph{two-pointer sweep over both arrivals and departures}.
\end{patternbox}

\subsection*{Understanding the problem carefully}

At any instant in time, the number of platforms required is exactly the
number of trains whose interval $[\textit{arr}[i], \textit{dep}[i]]$
currently contains that instant -- i.e., the number of trains
\emph{simultaneously} present at the station. The answer to the whole
problem is simply the \textbf{maximum}, over all instants in time, of this
simultaneous count. So the task reduces to efficiently finding the peak
number of overlapping intervals.

A direct way to track this is to imagine sweeping a clock forward through
time. Every arrival event increases the number of trains currently at the
station by one; every departure event decreases it by one. If we process
all arrival and departure events in time order, the running count after
each event tells us exactly how many platforms are occupied at that moment,
and the maximum value this running count ever reaches is our answer.

\begin{ideabox}[Greedy Choice]
Sort the arrival times and the departure times \emph{independently} (we do
not need to keep them paired to the same train -- only the multiset of
arrival instants and the multiset of departure instants matter for counting
overlaps). Use two pointers, one walking through sorted arrivals and one
through sorted departures. At each step, compare the next unprocessed
arrival to the next unprocessed departure: if the arrival happens at or
before that departure, a new train has shown up while the previous one has
not yet left, so increment the platform count needed and advance the
arrival pointer; otherwise, a train has left before this conflict arises, so
decrement the platform count and advance the departure pointer. Track the
maximum platform count seen at any point.
\end{ideabox}

\subsection*{Why sorting arrivals and departures independently is valid}

It may seem surprising that we are allowed to sort $\textit{arr}$ and
$\textit{dep}$ separately, losing the information of which arrival belongs
to which train. But the quantity we are computing -- the number of
simultaneously occupied platforms at any instant -- depends only on
\emph{how many} arrival events have occurred and \emph{how many} departure
events have occurred up to that instant, not on which specific train each
event belongs to. The running difference (arrivals processed so far minus
departures processed so far) is precisely the number of trains currently at
the station, regardless of how we have paired up arrival times with
departure times internally. This is what justifies treating $\textit{arr}$
and $\textit{dep}$ as two independent sorted streams of timestamps and
sweeping through them with two pointers, taking the running difference's
maximum value.

\subsection*{Algorithm}

\begin{enumerate}
  \item Sort $\textit{arr}$ in increasing order; sort $\textit{dep}$ in
        increasing order.
  \item Initialize $i \gets 0$, $j \gets 0$, $\textit{platforms} \gets 0$,
        $\textit{maxPlatforms} \gets 0$.
  \item While $i < n$ and $j < n$:
  \begin{enumerate}
    \item If $\textit{arr}[i] \le \textit{dep}[j]$: a new train has
          arrived while the earliest still-pending train has not yet left.
          $\textit{platforms} \gets \textit{platforms}+1$;
          $\textit{maxPlatforms} \gets \max(\textit{maxPlatforms},
          \textit{platforms})$; $i \gets i+1$.
    \item Otherwise: a train has departed.
          $\textit{platforms} \gets \textit{platforms}-1$; $j \gets j+1$.
  \end{enumerate}
  \item Return \textit{maxPlatforms}.
\end{enumerate}

\begin{algorithm}[H]
\caption{Minimum Platforms}
\begin{algorithmic}[1]
\Procedure{MinPlatforms}{$\textit{arr}, \textit{dep}$}
  \State sort $\textit{arr}$ in increasing order
  \State sort $\textit{dep}$ in increasing order
  \State $i \gets 0$, $j \gets 0$, $\textit{platforms} \gets 0$, $\textit{maxPlatforms} \gets 0$
  \While{$i < n$ \textbf{and} $j < n$}
    \If{$\textit{arr}[i] \le \textit{dep}[j]$}
      \State $\textit{platforms} \gets \textit{platforms} + 1$
      \State $\textit{maxPlatforms} \gets \max(\textit{maxPlatforms}, \textit{platforms})$
      \State $i \gets i + 1$
    \Else
      \State $\textit{platforms} \gets \textit{platforms} - 1$
      \State $j \gets j + 1$
    \EndIf
  \EndWhile
  \State \Return \textit{maxPlatforms}
\EndProcedure
\end{algorithmic}
\end{algorithm}

\subsection*{Dry run}

\dryrun{Take $\textit{arr} = [900, 940, 950, 1100, 1500, 1800]$ and
$\textit{dep} = [910, 1200, 1120, 1130, 1900, 2000]$.}

Sorted: $\textit{arr} = [900, 940, 950, 1100, 1500, 1800]$ (already sorted),
$\textit{dep} = [910, 1120, 1130, 1200, 1900, 2000]$ (sorted).

\begin{itemize}
  \item $\textit{arr}[0]=900 \le \textit{dep}[0]=910$. $\textit{platforms}=1$,
        $\textit{maxPlatforms}=1$, $i=1$.
  \item $\textit{arr}[1]=940 \le \textit{dep}[0]=910$? No ($940 > 910$).
        $\textit{platforms}=0$, $j=1$.
  \item $\textit{arr}[1]=940 \le \textit{dep}[1]=1120$. $\textit{platforms}=1$,
        $\textit{maxPlatforms}=1$, $i=2$.
  \item $\textit{arr}[2]=950 \le \textit{dep}[1]=1120$. $\textit{platforms}=2$,
        $\textit{maxPlatforms}=2$, $i=3$.
  \item $\textit{arr}[3]=1100 \le \textit{dep}[1]=1120$. $\textit{platforms}=3$,
        $\textit{maxPlatforms}=3$, $i=4$.
  \item $\textit{arr}[4]=1500 \le \textit{dep}[1]=1120$? No.
        $\textit{platforms}=2$, $j=2$.
  \item $\textit{arr}[4]=1500 \le \textit{dep}[2]=1130$? No.
        $\textit{platforms}=1$, $j=3$.
  \item $\textit{arr}[4]=1500 \le \textit{dep}[3]=1200$? No.
        $\textit{platforms}=0$, $j=4$.
  \item $\textit{arr}[4]=1500 \le \textit{dep}[4]=1900$. $\textit{platforms}=1$,
        $\textit{maxPlatforms}$ stays $3$, $i=5$.
  \item $\textit{arr}[5]=1800 \le \textit{dep}[4]=1900$. $\textit{platforms}=2$,
        $\textit{maxPlatforms}$ stays $3$, $i=6$.
  \item $i = n = 6$, loop ends.
\end{itemize}

The minimum number of platforms required is $\boxed{3}$, reached around
time $1100$--$1120$, when three trains are simultaneously at the station.

\subsection*{C++ implementation}

\begin{lstlisting}
#include <bits/stdc++.h>
using namespace std;

int minPlatforms(vector<int>& arr, vector<int>& dep) {
    int n = arr.size();
    sort(arr.begin(), arr.end());
    sort(dep.begin(), dep.end());

    int i = 0, j = 0;
    int platforms = 0, maxPlatforms = 0;

    while (i < n && j < n) {
        if (arr[i] <= dep[j]) {
            platforms++;
            maxPlatforms = max(maxPlatforms, platforms);
            i++;
        } else {
            platforms--;
            j++;
        }
    }
    return maxPlatforms;
}

int main() {
    vector<int> arr = {900, 940, 950, 1100, 1500, 1800};
    vector<int> dep = {910, 1200, 1120, 1130, 1900, 2000};

    cout << "Minimum platforms required: "
         << minPlatforms(arr, dep) << endl;
    return 0;
}
\end{lstlisting}

\begin{complexitybox}
\textbf{Time Complexity.} Sorting both arrays costs $O(n \log n)$. The
two-pointer sweep afterward visits each arrival and each departure exactly
once, costing $O(n)$. The overall time complexity is
\[
  O(n \log n).
\]
\textbf{Space Complexity.} Sorting can be performed in place (or with
$O(n)$ auxiliary space depending on the sorting routine), and the sweep
itself uses only a handful of scalar variables, so the auxiliary space
complexity is $O(1)$ beyond the input arrays.
\end{complexitybox}

\begin{takeawaybox}
When a problem asks for the minimum number of identical resources needed so
that \emph{every} request can be served without waiting, think in terms of
the maximum simultaneous overlap of intervals, computed by sweeping sorted
arrival and departure timestamps with a running counter. This is a different
flavor of greedy from activity selection: there, sorting by finish time let
us discard conflicting requests one at a time; here, sorting both endpoints
lets us measure the peak load directly, with nothing discarded.
\end{takeawaybox}
